% =============================================================================
% Chapter 1: Executive Summary
% Input-ready fragment for the lean toy-model research proposal.
% This file intentionally contains no preamble or document environment.
% =============================================================================

\chapter{Executive Summary}
\label{ch:executive-summary}

\section{The model and its central test}

This proposal defines a falsifiable research program for a classical toy analog.
It does not claim that the real universe is literally constituted in this way or
assume that its open targets have already been derived. The postulated arena has
preferred time and four Euclidean spatial dimensions. One globally conserved
compressible medium occupies two material states: an ordered, finite-thickness,
shear-supporting brane-like state and a de-structured bulk-like state. The brane
and bulk are therefore not separate substances; all transport, conversion, return,
radiation, and reservoir response must ultimately close within the same medium.

Light is sought as the lowest gapless transverse branch guided by the ordered
brane, with two freely selectable polarizations. This conditional linear-light
sector is distinct from the stronger nonlinear hypothesis of a finite-energy,
three-dimensionally localized classical photon packet. Particles are hypothesized
to be finite nonlinear throats through the brane. A trapped transverse standing
mode may help hold a throat open, but its support energy is structural and is not
identified automatically with observed particle mass. The side of the brane toward
which a throat is oriented supplies the proposed electric sign, while particle
species and a complete antiparticle map require additional internal, support-mode,
topological, or composite structure.

The other observable sectors are intended to arise from the response of the same
medium. Localized one-way throat drainage, throat-region pressure, stress, and
order conversion, the separate distributed porous-brane reordering encoded by
the positive bulk-to-brane contribution to \(S_{\mathrm{leak}}\), and reservoir
response are proposed as the source-side material basis for a calibrated gravity
observable. Motion of an oriented throat is expected to produce a sign-dependent
magnetic response, while acceleration must be tested for emission into both
transverse light polarizations. Inertia is attributed to reconfiguration of the
complete dressed object: its throat geometry, trapped mode, near fields, flow,
conversion region, and surrounding medium. Active gravitational strength,
passive gravitational response, and inertial resistance are kept distinct until
their signs and low-energy universality are derived.

The central question is whether one mathematically coherent parent medium can
support all mandatory sectors in compatible regimes with one field content, one
dynamical branch, one coefficient set, and one closed treatment of boundary and
reservoir data. If \(\Lambda_S\) is the admissible region left by sector \(S\),
the decisive requirement is

\[
\Lambda_{\mathrm{common}}
=
\bigcap_{S\in\mathcal S_{\mathrm{required}}}\Lambda_S
\neq\varnothing.
\]

A successful light point, electric point, and gravity point selected from
incompatible branches or parameter regions would not satisfy this test.

\section{The far-field-first research program}

Because the final parent equations and nonlinear throats are not yet closed, the
program begins with far-field capability requirements rather than allowing an
unresolved particle core to supply whatever source a later calculation needs. For
each sector, the proposal first states what a brane observer must see, identifies
the proposed material mechanism, specifies the static and dynamic calculations
needed to test it, and extracts what the result requires or forbids in the common
medium. Requirement discovery works backward from observable behavior;
mathematical closure later works forward from one frozen parent system to its
background, spectrum, sources, and nonlinear solutions.

Light leads because it directly tests the host medium. The first track asks
whether a stable finite slab supports a guided, gapless transverse branch with two
polarizations, acceptable confinement across the normal direction, and controlled
dispersion, leakage, and mixing. The second asks whether finite-thickness
dispersion and a slaved reflection-even material dressing can maintain a localized
photon-packet candidate without collapse, strong generic interaction, or
unacceptable radiation. Failure of the packet track would not erase a valid
conditional classical light-wave sector, but it would reject the intended
self-bound free-photon interpretation.

Gravity asks whether localized one-way throat flow, stress, and order conversion,
together with the separate distributed porous-brane \(S_{\mathrm{leak}}\) return
and the resulting reservoir response, can produce a probe-independent field with
an attractive inverse-square local regime, orientation-even source behavior, one
frozen calibration, acceptable dynamics, and calculable return-induced screening
or overshoot. Static electricity
asks whether the complete finite-slab response contains a healthy reflection-odd
carrier with leading \(k^2\) static stiffness, the required interaction sign,
sufficiently universal throat coupling, and a conserved oriented count or current.
Parity or gaplessness alone is not enough, and every additional long-range even or
odd mode must be included.

Magnetism and radiation then test whether electricity belongs to one dynamical
source hierarchy rather than remaining an isolated static analogy. The same
oriented source and coupling must govern the leading static electric response, the
velocity-order magnetic response, and acceleration-order emission into both
transverse light polarizations. Radiation into odd, even, longitudinal, bulk, and
order-conversion or reservoir branches must be calculated rather than omitted. The
far-field phase also establishes requirements for positive passive and inertial
response, retarded passivity, low-energy species universality, compatible
characteristic speeds, and acceptable leakage, emission, friction, and drag.

Each sector hands its required fields, operators, speeds, source conventions,
normalizations, conservation laws, permitted departures, and failure conditions to
a common Medium Capability Matrix. This exposes shared quantities and conflicts
before sector-specific repair can hide them.

\section{From sector requirements to one common medium}

The sector outputs are normalized and intersected. A mandatory failure cannot be
compensated by strong performance in another sector, and the union of
independently successful parameter points is not a common solution. The surviving
region must be reported honestly as empty, isolated, fragile, disconnected, or
robust. The compatibility tests are potential falsification gates, not
established contradictions.

If a nonempty region remains, the project will freeze Candidate Parent System V1.
The freeze must specify the arena, field inventory, conservative, relaxational, or
mixed dynamical branch, equations and constraints, operator basis and derivative
order, symmetries, independent parameters, boundary and reservoir data, source and
observable conventions, and terms forbidden from later sector-specific insertion.
A coefficient selected to localize a photon must also propagate through the slab
spectrum, electric response, throat support, radiation, and drag. A coupling fixed
by static electricity must retain its role in magnetism and acceleration
radiation. Any substantive revision creates a new candidate version and reopens
its actual downstream dependencies.

Candidate V1 must then produce a finite, selected, stable or sufficiently
metastable brane and its complete constrained spectrum. All far-field sectors must
be rederived from that same background and common parameter region. Only this
bottom-up rederivation can show that the capability contracts are realized rather
than merely compatible on paper. The proposal supplies the top-down contract; the
derivation ledger supplies the executable evidence, assumptions, frozen inputs,
and dependency record used to decide it.

\section{Beyond the far field and meaning of success or failure}

Only a candidate that survives the common-background and far-field rederivation
advances to the nonlinear particle program. The first task is to continue a
localized transverse bound-mode seed or acceptably long-lived resonance on a
provisional throat geometry into a self-consistent supported throat whose shape,
support stress, transport, conversion, and reservoir loading agree. Its stability,
source profiles, and couplings to the complete background spectrum must then be
established.

Both orientations of each species must be solved from the same equations. Electric
orientation, species identity, and antiparticle identity remain separate: a
complete antiparticle candidate may require transformation of internal phase,
circulation, chirality, or topology in addition to geometric reflection. Moving
and accelerating throat families must determine inertia, passive response,
magnetism, radiation, drag, and any low-energy universality among active, passive,
and inertial roles. Two-throat calculations must continue from the far field
through the finite-thickness crossover and into full nonlinear core geometry.
Projected coincidence, cancellation of the leading odd electric disturbance,
four-dimensional core contact, and annihilation are distinct events.
Same-species throat--antithroat annihilation must therefore be distinguished
from possible neutral binding of oppositely charged objects with different
internal or composite structure.

Finally, the local throat and global medium must be solved as one feedback
problem. Localized one-way throat drainage and throat-region order conversion
alter the bulk reservoir; material does not return through the throat. Instead,
because the finite-thickness brane is porous, bulk material can reorganize into
the ordered brane state across it. For a projected brane observer, that
distributed bulk-to-brane return is the positive contribution to the signed term
\(S_{\mathrm{leak}}\). It changes the brane and the asymptotic data seen by the
throat. Closure requires a fixed point with consistent material, momentum,
energy, and, where relevant, entropy accounting. A local throat solution is not
complete if its own sustained outflow would invalidate the boundary data used to
obtain it.

The possible outcomes are deliberately layered. A valid linear-light sector is a
result even if the localized photon hypothesis fails. Individually viable
far-field sectors are useful but do not establish one common medium. Failure of
Candidate V1 rejects that candidate, not automatically every allowed realization
of the broader architecture. The architecture fails only if its mandatory
capabilities prove mutually incompatible throughout the admitted frozen theory
class or require prohibited retuning or an unsupported mechanism.

Full success requires one controlled parent system to survive the common
far-field intersection, the stable-brane and complete-spectrum tests, the
supported-throat program, moving and interacting solutions, and global reservoir
closure. Even that result would establish a
classical analog architecture only within its stated scope, not exact general
relativity, exact electromagnetism, quantum theory, or a derivation of the real
universe. An empty common region, unavoidable wrong-sign or extra long-range mode,
excessive leakage or drag, failed photon-packet or throat branch, unacceptable
species dependence, or unclosed reservoir cycle is an equally legitimate result
because it identifies where the proposed architecture fails.
